%% parts/p02/first-order-normal-forms-supplement.tex --- retained additions
\section*{Editorial normal-form additions}
\addcontentsline{toc}{section}{Editorial normal-form additions}
\noindent\textit{Reader note.} The exact inherited chapter supplies the governing
normal-form definitions, results, and proofs.  The following exact historical
slices retain only the quantifier-law table, normal-form comparison, and
clausal-pipeline figure.
They remain separately attributed and do not replace or alter the inherited
source.

\begin{hbtheorem}{thm:p02-fol-quantifier-equivalences}{cited=Enderton2001}{Quantifier equivalence laws}
Every row of Table~\ref{tab:p02-quantifier-equivalences} is valid in the
nonempty-domain first-order semantics.  A row carrying
\(x\notin\FV(\chi)\) may be used only under that displayed side condition.
\end{hbtheorem}

\begingroup\small\setlength{\tabcolsep}{4pt}
\begin{longtable}{@{}p{.32in}p{2.55in}p{1.88in}@{}}
\caption{Quantifier equivalence laws with side conditions.}
\label{tab:p02-quantifier-equivalences}\\
\toprule
Row & Equivalence & Side condition \\ \midrule
\endfirsthead
\toprule Row & Equivalence & Side condition \\ \midrule
\endhead
Q1 & \(\neg\forall x\,\varphi\Leftrightarrow
       \exists x\,\neg\varphi\) & none \\
Q2 & \(\neg\exists x\,\varphi\Leftrightarrow
       \forall x\,\neg\varphi\) & none \\
Q3 & \((Qx\,\varphi)\land\chi\Leftrightarrow
       Qx(\varphi\land\chi)\), \(Q\in\{\forall,\exists\}\) &
       \(x\notin\FV(\chi)\) \\
Q4 & \((Qx\,\varphi)\lor\chi\Leftrightarrow
       Qx(\varphi\lor\chi)\), \(Q\in\{\forall,\exists\}\) &
       \(x\notin\FV(\chi)\) \\
Q5 & \((\forall x\,\varphi)\to\chi\Leftrightarrow
       \exists x(\varphi\to\chi)\) & \(x\notin\FV(\chi)\) \\
Q6 & \((\exists x\,\varphi)\to\chi\Leftrightarrow
       \forall x(\varphi\to\chi)\) & \(x\notin\FV(\chi)\) \\
Q7 & \(\chi\to(\forall x\,\varphi)\Leftrightarrow
       \forall x(\chi\to\varphi)\) & \(x\notin\FV(\chi)\) \\
Q8 & \(\chi\to(\exists x\,\varphi)\Leftrightarrow
       \exists x(\chi\to\varphi)\) & \(x\notin\FV(\chi)\) \\
Q9 & \(\forall x(\varphi\land\psi)\Leftrightarrow
       (\forall x\varphi)\land(\forall x\psi)\) & none \\
Q10 & \(\exists x(\varphi\lor\psi)\Leftrightarrow
         (\exists x\varphi)\lor(\exists x\psi)\) & none \\
Q11 & \(\forall x\,\varphi\Leftrightarrow\varphi\) &
        \(x\notin\FV(\varphi)\) \\
Q12 & \(\exists x\,\varphi\Leftrightarrow\varphi\) &
        \(x\notin\FV(\varphi)\) \\
\bottomrule
\end{longtable}
\endgroup

The converses of the familiar one-way distribution laws omitted from the
table are generally invalid; for example, two existential conjuncts need not
share a witness.  The side conditions in Q3--Q8 are part of the law, not
typographical annotations.


\section{The clausal-form pipeline}
\label{sec:p02-fol-clausal-pipeline}

\begin{hbtheorem}{thm:p02-fol-clausal-form}{cited=BaazEglyLeitsch2001}{Clausal-form theorem}
For every first-order sentence \(\sigma\), there is a finite set \(S\) of
universally read clauses in a finite expansion of its signature such that,
assuming AC, \(\sigma\) is satisfiable iff \(S\) is satisfiable.  The reverse
direction from an \(S\)-model to a model of \(\sigma\) uses only reducts and
does not require choice.  The transformation is the composition in
Table~\ref{tab:p02-fol-normal-form-pipeline}; only the Skolem step weakens
same-signature logical equivalence to expansion--reduct equisatisfiability.
\end{hbtheorem}

\begin{landscape}
\begingroup\scriptsize\setlength{\tabcolsep}{3pt}
\begin{longtable}{@{}p{.30in}p{1.05in}p{1.65in}p{1.90in}p{2.00in}@{}}
\caption{First-order normal-form comparison and clausal pipeline.}
\label{tab:p02-fol-normal-form-pipeline}\\
\toprule
Step & Output & Existence and scope & Preservation claim & Worst-case size and new data \\ \midrule
\endfirsthead
\toprule Step & Output & Existence and scope & Preservation claim & Worst-case size and new data \\ \midrule
\endhead
1 & NNF & Exists for every formula in the same signature. & Formula
equivalence; exactly the same free and all-variable sets. & Linear in the
formula tree; no new symbols. \\
2 & Rectified NNF & Exists because the variable supply is countably infinite;
fresh binders are chosen outside a finite forbidden set. & Alpha-equivalence,
hence formula equivalence; same constructor tree. & Linear apart from the
encoding of fresh variable names; no new nonlogical symbols. \\
3 & Prenex NNF & Exists for every rectified NNF formula by the side-conditioned
laws in Table~\ref{tab:p02-quantifier-equivalences}. & Formula equivalence;
same free and all-variable sets and the same quantifier occurrences. & Linear
in the formula tree; only the prefix is rearranged. \\
4 & Universal Skolem sentence & Exists for every rectified prenex sentence in a
finite signature expansion; AC supplies witness functions in the forward
model-expansion direction. & Expansion--reduct equisatisfiability under AC,
not same-signature formula equivalence. & At most quadratic under literal term
printing because a new Skolem term may list all preceding universal variables;
new function symbols are introduced. \\
5 & Matrix in CNF & Exists for every quantifier-free matrix by propositional
distribution. & Formula equivalence in the expanded signature. & Direct
distribution can be exponential; Tseitin conversion is only equisatisfiable
and introduces symbols (see Part~2.1). \\
6 & Matrix in DNF & Exists for every quantifier-free matrix by dual
propositional distribution. & Formula equivalence in the expanded signature. &
Direct distribution can be exponential; DNF is a comparison form, not the
input used by resolution. \\
7 & Clause set & Exists from a universal sentence whose matrix is in CNF. & The
set and the universal closures of its clauses have exactly the same models as
that universal sentence. & Linear once the CNF matrix is present; erases the
written universal prefix and records a finite literal set per conjunct. \\
\bottomrule
\end{longtable}
\endgroup
\end{landscape}

\begin{figure}[htbp]
\centering
\begin{tikzpicture}[
  font=\small,
  node distance=8mm,
  box/.style={draw,rounded corners,align=center,inner sep=4pt,
              minimum width=1.55in},
  >=Stealth]
  \node[box] (in) {sentence\\\(\sigma\)};
  \node[box,below=of in] (nnf) {NNF};
  \node[box,below=of nnf] (rect) {rectified NNF};
  \node[box,below=of rect] (pren) {rectified prenex NNF};
  \node[box,below=of pren] (sk) {universal Skolem sentence};
  \node[box,below=of sk] (cnf) {universal sentence\\with CNF matrix};
  \node[box,below=of cnf] (clauses) {clause set \(S\)};
  \draw[->] (in) -- node[right=8mm,align=left,font=\footnotesize]
    {formula equivalence} (nnf);
  \draw[->] (nnf) -- node[right=8mm,align=left,font=\footnotesize]
    {alpha-equivalence} (rect);
  \draw[->] (rect) -- node[right=8mm,align=left,font=\footnotesize]
    {formula equivalence} (pren);
  \draw[->] (pren) -- node[right=8mm,align=left,font=\footnotesize]
    {equisatisfiable\\expanded signature} (sk);
  \draw[->] (sk) -- node[right=8mm,align=left,font=\footnotesize]
    {matrix CNF\\formula equivalence} (cnf);
  \draw[->] (cnf) -- node[right=8mm,align=left,font=\footnotesize]
    {universal-closure convention} (clauses);
\end{tikzpicture}
\caption{Normal-form pipeline.  The Skolem arrow records
expansion--reduct equisatisfiability rather than same-signature formula
equivalence; the last arrow relates a sentence to a clause set through the
universal closures of its members, not through formula-string equivalence.}
\label{fig:p02-fol-normal-form-pipeline}
\end{figure}

This theorem composes the preceding transformations at their stated proof
statuses.  Resolution
objects and their universal-closure reading are defined in the next chapter;
the table does not identify a finite clause set with a single formula string.
For the quantifier equivalences and the clause-conversion pipeline,
respectively, see \cite{Enderton2001,BaazEglyLeitsch2001}.
